%------------------------------------------------------------------------------%
% Luis A. D'Afonseca
%
%------------------------------------------------------------------------------%
\documentclass[a4paper,12pt,fleqn]{article}

\usepackage{style_avaliacao}

\ShowAnswers

\newcommand{\D}[2]{\frac{\partial {#1}}{\partial {#2}}}
\newcommand{\DS}[2]{\frac{\partial^2 {#1}}{\partial {#2}^2}}
\newcommand{\DM}[3]{\frac{\partial^2 {#1}}{\partial {#2} \partial {#3}}}

\newcommand{\vetor}[2]{\left(\begin{array}{c} #1 \\ #2 \end{array}\right)}

\newcommand{\veci}{\bm{i}}
\newcommand{\vecj}{\bm{j}}
\newcommand{\veck}{\bm{k}}

\newcommand{\comentario}[1]{{\color{magenta} #1 \par\vspace{\baselineskip}\par}}

%------------------------------------------------------------------------------%
\begin{document}

\makeHeader{CFVVI}{Prova 3 -- Turma 6}{09/09/2024}

%\comentario{
%  Nos exercícios sobre números complexos:
%  \begin{itemize}
%    \item Como esta é uma disciplina de cálculo, exigi que as respostas sejam sempre
%    em radianos, mas os alunos podem fazer os passos intermediários em graus.
%
%    \item No cálculo das raizes o aluno pode deixar a resposta final na forma polar,
%    não é necessário retornar para a forma canônica.
%  \end{itemize}
%}

% 01
%------------------------------------------------------------------------------%
\exercise{50}
Encontre o ponto do plano \(x+2y+3z=13\) mais próximo do ponto \((1, 1, 1)\).

\begin{answer}
  Queremos minimizar a distância
  \[
    d(x, y, z) = \sqrt{(x-1)^2 + (y-1)^2 + (z-1)^2}
  \]
  Como a função raiz quadrada é crescente podemos minimizar a função
  \[
    f(x, y, z) = (x-1)^2 + (y-1)^2 + (z-1)^2
  \]
  suas derivadas parciais são
  \[
    f_x(x, y, z) = 2(x-1)
    \qquad
    f_y(x, y, z) = 2(y-1)
    \qquad
    f_z(x, y, z) = 2(z-1)
  \]
  As derivadas parciais da função \(g(x, y, z)=x+2y+3z\) são
  \[
    g_x(x, y, z) = 1
    \qquad
    g_y(x, y, z) = 2
    \qquad
    g_z(x, y, z) = 3
  \]
  Aplicando os Multiplicadores de Lagrange, precisamos resolver o sistema
  \[
  \begin{cases}
    2(x-1) =  \lambda \\
    2(y-1) = 2\lambda \\
    2(z-1) = 3\lambda \\
    x + 2y + 3z = 13
  \end{cases}
  \]
  Das três primeiras equações obtemos
  \[
    2(x-1) = \lambda
    \quad\Rightarrow\quad
    x = \frac{\lambda}{2} + 1
  \]
  \[
    2(y-1) = 2\lambda
    \quad\Rightarrow\quad
    y = \lambda + 1
  \]
  \[
    2(z-1) = 3\lambda
    \quad\Rightarrow\quad
    z = \frac{3\lambda}{2} + 1
  \]
  Substituindo na última equação
  \begin{align*}
    x + 2y + 3z & = 13
    \\[2mm]
    \left(\frac{\lambda}{2} + 1\right)
    + 2 \left(\lambda + 1\right)
    + 3 \left(\frac{3\lambda}{2} + 1\right) & = 13
    \\[2mm]
    \frac{\lambda}{2} + 1
    + 2\lambda + 2
    + \frac{9\lambda}{2} + 3 & = 13
    \\[2mm]
    \lambda + 2
    + 4\lambda + 4
    + 9\lambda + 6 & = 13\times 2
    \\[2mm]
    14\lambda + 12 & = 26
    \\[2mm]
    14\lambda & = 14
    \\[2mm]
    \lambda & = 1
  \end{align*}
  Portanto
  \[
    x = \frac{\lambda}{2} + 1 = \frac{3}{2}
  \]
  \[
    y = \lambda + 1 = 2
  \]
  \[
    z = \frac{3\lambda}{2} + 1 = \frac{5}{2}
  \]
  O ponto mais próximo é o ponto \(\left(\frac{3}{2}, 2, \frac{5}{2}\right)\)
  \clearpage
\end{answer}

% 02
%------------------------------------------------------------------------------%
\exercise{25}
Encontre as raízes do polinômio \(p(x) = x^3 -6x^2 + 10x\)

\begin{answer}
  Queremos resolver a equação
  \[
    x^3 -6x^2 + 10x = 0
  \]
  que podemos escrever como
  \[
    x\left(x^2 -6x + 10\right) = 0
  \]
  Uma das raízes é \(x=0\), usamos Bhaskara para encontrar as demais
  \[
    \Delta
    = b^2 - 4ac
    = (-6)^2 - 4 \times 1 \times 10
    = 36 - 40
    = -4
  \]
  \[
    \sqrt{\Delta}
    = \sqrt{-4}
    = i\sqrt{4}
    = 2i
  \]
  Temos então
  \[
    x
    = \frac{-b\pm\sqrt{\Delta}}{2a}
    = \frac{-(-6)\pm 2i}{2}
    = 3\pm i
  \]
  Portanto as raízes de $p$ são
  \[
    x_1 = 0 \qquad\qquad
    x_2 = 3-i \qquad\qquad
    x_3 = 3+i
  \]
  \clearpage
\end{answer}

% 03
%------------------------------------------------------------------------------%
\exercise{25}
Encontre as raízes cúbicas complexas do número \( z = -27 \).

\begin{answer}
  Escrevemos \( z = -27 \) na forma trigonométrica
  \[
  z = 27 \left[ \cos\left(\pi\right) + i \sin\left(\pi\right) \right]
  \]
  As raízes cúbicas são
  \begin{align*}
    u_k
    & = \sqrt[3]{27} \left[
    \cos\left(\frac{\pi+2k\pi}{3}\right)
    + i \sin\left(\frac{\pi+2k\pi}{3}\right)
    \right]
    & k = 0, 1, 2    \\[2mm]
    & = 3 \left[
    \cos\left(\frac{\pi+2k\pi}{3}\right)
    + i \sin\left(\frac{\pi+2k\pi}{3}\right)
    \right]
  \end{align*}
  Para cada valor de $k$
  \begin{align*}
    u_0
    & = 3 \left[
    \cos\left( \frac{\pi+2\times 0\pi}{3} \right)
    + i \sin\left( \frac{\pi+2\times 0\pi}{3} \right)
    \right] \\[2mm]
    & = 3 \left[
    \cos\left( \frac{\pi}{3} \right)
    + i \sin\left( \frac{\pi}{3} \right)
    \right] \\[2mm]
    & = 3 \left[
    \cos\left( \ang{60} \right)
    + i \sin\left( \ang{60} \right)
    \right] \\[2mm]
    & = 3 \left[
    \frac{1}{2}
    + i \frac{\sqrt{3}}{2}
    \right] \\[2mm]
    & =  \frac{3}{2} + \frac{3\sqrt{3}}{2} i
  \end{align*}
  \begin{align*}
    u_1
    & = 3 \left[
    \cos\left( \frac{\pi+2\times 1\pi}{3} \right)
    + i \sin\left( \frac{\pi+2\times 1\pi}{3} \right)
    \right] \\[2mm]
    & = 3 \left[
    \cos\left( \pi \right)
    + i \sin\left( \pi \right)
    \right] \\[2mm]
    & = 3 \left[ -1 + i0 \right] \\[2mm]
    & = -3
  \end{align*}

  \begin{align*}
    u_2
    & = 3 \left[
    \cos\left( \frac{\pi+2\times 2\pi}{3} \right)
    + i \sin\left( \frac{\pi+2\times 2\pi}{3} \right)
    \right] \\[2mm]
    & = 3 \left[
    \cos\left( \frac{5\pi}{3} \right)
    + i \sin\left( \frac{5\pi}{3} \right)
    \right] \\[2mm]
    & = 3 \left[
    \cos\left( \ang{300} \right)
    + i \sin\left( \ang{300} \right)
    \right] \\[2mm]
    & = 3 \left[
    \frac{1}{2}
    - i \frac{\sqrt{3}}{2}
    \right] \\[2mm]
    & = \frac{3}{2} - \frac{3\sqrt{3}}{2}i
  \end{align*}
\end{answer}

%------------------------------------------------------------------------------%
\end{document}
%------------------------------------------------------------------------------%
